\documentclass[a4paper,11pt]{article}
\usepackage{jinstpub}

% Packages
\usepackage{graphicx}
\usepackage{amsmath,amssymb}
\usepackage{booktabs}

% Title and author
\title{Design and Simulation of a Nonlinear Kicker Magnet for Booster-to-Storage Ring Injection}

\author[a]{Chong Shik Park}

\affiliation[a]{Department of Accelerator Science, Korea University, Sejong, 30019 South Korea}

\emailAdd{kuphy@korea.ac.kr}

\abstract{
We present the modeling, simulation, and optimization of a Booster-to-Storage Ring (BTS) beam transport line incorporating a nonlinear kicker magnet for injection into a storage ring. This study focuses on the design of a nonlinear kicker that enables efficient injection while minimizing emittance growth and preserving beam quality. Detailed field calculations reveal the spatial nonlinearity and asymmetry of the kicker field, which is then incorporated into beam dynamics simulations. The BTS lattice is optimized to match the Twiss parameters at the injection point while maintaining transport efficiency. Phase space analysis before and after kicker passage shows significant deformation, highlighting the importance of nonlinear field modeling in injection design. This work serves as a foundation for implementing advanced kicker technologies in modern light source facilities.
}

\keywords{Beam dynamics, Injection, Synchrotron radiation, Nonlinear magnets, Accelerator simulation}

\begin{document}
\maketitle
\flushbottom

\section{Introduction}

Efficient beam injection from a booster synchrotron to a storage ring is a critical aspect of modern synchrotron light sources. Traditional kicker magnets used for injection often introduce unwanted emittance growth and beam quality degradation due to their limited field profiles and space constraints. To address these challenges, novel nonlinear kicker magnets have been proposed, offering greater control over the injected beam trajectory while potentially reducing the perturbation to the stored beam.

This study presents a design and simulation-based analysis of a nonlinear kicker magnet incorporated into a Booster-to-Storage Ring (BTS) transfer line. We aim to evaluate the performance of the nonlinear kicker in the context of beam transport, optical matching, and phase space preservation. By coupling detailed magnetic field modeling with Twiss parameter optimization and beam tracking simulations, we explore the feasibility and advantages of such a system.

\section{Design of the BTS Beamline}

The BTS beamline connects the booster synchrotron to the storage ring and must ensure proper matching of the beam's optical parameters at the injection point. The design goal is to preserve beam quality while adapting the beam envelope to the requirements of the storage ring.

\subsection{Twiss Parameter Matching}

The initial Twiss parameters at the booster extraction point are given by:
\begin{equation}
    (\beta_x, \alpha_x, \beta_y, \alpha_y)_{\text{in}} = (4.25\ \mathrm{m}, -0.81,\ 2.77\ \mathrm{m}, -0.27).
\end{equation}

Through iterative lattice tuning and quadrupole strength adjustment, the final Twiss parameters at the end of the transfer line are matched to:
\begin{equation}
    (\beta_x, \alpha_x, \beta_y, \alpha_y)_{\text{out}} = (4.19\ \mathrm{m}, -0.78,\ 2.65\ \mathrm{m}, -0.25).
\end{equation}

These values show that the beamline preserves the optical functions with minimal mismatch, supporting stable injection conditions.

\subsection{Lattice Configuration}

The BTS lattice is composed of a sequence of quadrupole magnets, drifts, and diagnostic elements. The quadrupole strengths were optimized to minimize the mismatch factor $\mathcal{M}$, defined as:
\begin{equation}
    \mathcal{M} = \frac{1}{2} \left( \frac{\beta_{\text{out}}}{\beta_{\text{target}}} + \frac{\beta_{\text{target}}}{\beta_{\text{out}}} - 2 \right) + (\alpha_{\text{out}} - \alpha_{\text{target}})^2.
\end{equation}

Tracking and optics matching were carried out using custom Python tools developed within the SynapticTrack framework.

\section{Nonlinear Kicker Magnet}

\subsection{Motivation and Concept}

Conventional kicker magnets produce uniform transverse fields that impart linear kicks to the beam. However, for off-axis injection schemes or selective beam manipulation, nonlinear field profiles can offer tailored control. A nonlinear kicker magnet is designed to produce a position-dependent transverse field, enabling stronger deflections at larger offsets and minimal perturbation near the axis.

\subsection{Magnetic Field Modeling}

The magnetic field of the nonlinear kicker was modeled using a numerical field map. Figure~\ref{fig:field_map} shows the simulated vertical magnetic field $B_y$ in the $x$--$z$ plane. The field is localized and exhibits strong transverse gradient nonlinearity near the magnet center.

\begin{figure}[htbp]
    \centering
    \includegraphics[width=0.6\textwidth]{field_map.png}
    \caption{Magnetic field distribution $B_y(x, z)$ of the nonlinear kicker magnet. The field is anti-symmetric in $x$ and decays rapidly along the longitudinal direction $z$.}
    \label{fig:field_map}
\end{figure}

\subsection{Field Gradient Profile}

To better understand the nonlinear characteristics, the transverse field gradient $\partial B_y/\partial x$ was extracted at various longitudinal positions. Figure~\ref{fig:field_grad} displays the 3D plot of field gradients, confirming the cubic-like dependence on transverse position $x$ near the magnet center.

\begin{figure}[htbp]
    \centering
    \includegraphics[width=0.65\textwidth]{field_gradient_3d.png}
    \caption{Transverse field gradient $\partial B_y / \partial x$ versus $x$ and $z$. The kicker field exhibits strong nonlinear dependence in $x$.}
    \label{fig:field_grad}
\end{figure}

\section{Beam Dynamics Simulations}

\subsection{Tracking Setup}

Particle tracking was performed through the BTS beamline and nonlinear kicker using a custom symplectic integrator. The field map of the kicker was interpolated to provide continuous transverse kicks during integration. The beam was initialized with Gaussian transverse distributions and matched Twiss parameters at the kicker entrance.

\subsection{Phase Space Evolution}

Figure~\ref{fig:phase_before_after} shows the phase space $(x, x')$ distribution before and after the kicker. The nonlinear field introduces clear distortions to the phase space, particularly at large amplitudes, as expected from the position-dependent kick.

\begin{figure}[htbp]
    \centering
    \includegraphics[width=0.49\textwidth]{phase_space_before.png}
    \includegraphics[width=0.49\textwidth]{phase_space_after.png}
    \caption{Horizontal phase space $(x, x')$ before (left) and after (right) the nonlinear kicker. The distorted distribution after the kicker reflects the nonlinear transverse kick.}
    \label{fig:phase_before_after}
\end{figure}

\subsection{Impact on Beam Emittance}

Despite the nonlinear phase space deformation, the projected emittance growth remains small due to careful beam optics matching and limited phase space area exposed to strong nonlinearity. Table~\ref{tab:emittance} summarizes the transverse emittances before and after the kicker.

\begin{table}[htbp]
    \centering
    \caption{RMS transverse emittance before and after nonlinear kicker.}
    \label{tab:emittance}
    \begin{tabular}{lcc}
        \toprule
        & Before Kicker & After Kicker \\
        \midrule
        $\varepsilon_x$ [mm·mrad] & 2.5 & 2.6 \\
        $\varepsilon_y$ [mm·mrad] & 2.4 & 2.4 \\
        \bottomrule
    \end{tabular}
\end{table}

\section{Results}

The simulation results validate the feasibility of incorporating a nonlinear kicker magnet into the BTS beamline. Key observations include:

\begin{itemize}
    \item \textbf{Field profile:} The kicker generates a highly nonlinear transverse field with strong gradients localized near the center.
    \item \textbf{Beam optics:} Twiss parameters are preserved with high accuracy across the BTS line, even with the nonlinear kicker.
    \item \textbf{Phase space effects:} The kicker introduces expected distortions in the transverse phase space but does not cause significant emittance growth.
\end{itemize}

The integration of nonlinear field effects into the beam dynamics simulation was essential to accurately capture the kick-induced distortions. This underscores the importance of detailed field modeling in kicker magnet design.

\acknowledgments

The author acknowledges the support of KAERI and valuable discussions with collaborators in the beam physics group. Computations were carried out using resources provided by the SynapticTrack simulation framework.

\bibliographystyle{JHEP}
\bibliography{nonlinear_kicker_refs}

\end{document}
